% Point to the main file
\documentclass[../../Main.tex]{subfiles}

% ========== Start the document ==========
\begin{document}

\section{Exam 4 Review}

\subsection{Binary and Base n}

Write 1577 in binary

\begin{align*}
    1577 &= 2 \cdot 788 + \yellowbox{1} \\
    788 &= 2 \cdot 394 + \yellowbox{0} \\
    394 &= 2 \cdot 197 + \yellowbox{0} \\
    197 &= 2 \cdot 98 + \yellowbox{1} \\
    98 &= 2 \cdot 49 + \yellowbox{0} \\
    49 &= 2 \cdot 24 + \yellowbox{1} \\
    24 &= 2 \cdot 12 + \yellowbox{0} \\
    12 &= 2 \cdot 6 + \yellowbox{0} \\
    6 &= 2 \cdot 3 + \yellowbox{0} \\
    3 &= 2 \cdot 1 + \yellowbox{1} \\
    1 &= 2 \cdot 0 + \yellowbox{1} \\
\end{align*}

Writing the bottom to top left to right...

$$\boxed{11000101001}$$

We could check this by just doing 

$$1577 = 2^{10} + 2^9 + \ldots $$

Let's try representing 1577 in base 8. Remember that each remainder we get is the coefficient to the corresponding power

\begin{align*}
    1577 &= 8 \cdot 197 + \yellowbox{1} \\
    197 &= 8 \cdot 24 + \yellowbox{5} \\
    24 &= 8 \cdot 3 + \yellowbox{0} \\
    3 &= 8 \cdot 0 + \yellowbox{3} \\
\end{align*}

Rewriting this, we get

\begin{align*}
    \underbrace{3}_{8^3}
    \underbrace{0}_{8^3}
    \underbrace{5}_{8^2}
    \underbrace{1}_{8^0}
\end{align*}

So 

$$1577 = 3 \cdot 8^3 + 5 \cdot 8^1 + 1 \cdot 8^0$$

Now let's do it with base 16

\begin{align*}
    1577 &= 16 \cdot 98 + \yellowbox{9} \\
    98 &= 16 \cdot 6 + \yellowbox{2} \\
    6 &= 16 \cdot 0 + \yellowbox{6} \\
\end{align*}

$$\boxed{1577 = 629_{16}}$$

$$1577 = 6 \cdot 16^2 + 2 \cdot 16^1 + 9 \cdot 16^0$$

For our last example, write the $AB0CD_{16}$ in base 10 (aka normal numbers)

Recall that

\begin{align*}
    A &= 10 \\
    B &= 11 \\
    &\vdots \\
\end{align*}

$AB0CD$ is just $10 \; 11 \; 0 \; 12 \; 13$

Since there's 5 "numbers" that means we start with $16^{5 - 1} = 16^4$

\begin{align*}
    AB0CD_{16} &= 10 \cdot 16^4 + 11 \cdot 16^3 + 0 \cdot 16^2 + 12 \cdot 16^1 + 13 \cdot 16^1 \\
    &= \boxed{700621} \tag{\text{just plug into calc}}
\end{align*}

\subsection{RSA and Encryption}

\subsubsection{FLT Encrytion}

For Encryption using FLT, basically

We have:

\begin{align*}
    M: &\; \text{Message} \\
    C: &\; \text{Encrypted Message} \\
    e: &\; \text{Encryption Key} \\
    d: &\; \text{Decryption Key} \\
    p: &\; \text{Prime mod} \\
\end{align*}

To \textbf{encrypt} a message $M$ to $C$

$$C \equiv M^e \pmod{p}$$

To \textbf{decrypt} $C$ back into $M$, we must first find $d$

$$ed \equiv 1 \pmod{p - 1}$$

Just find $e^{-1}$ using Euclidean and plug it into

$$M \equiv C^d \pmod{p}$$

\subsubsection{FLT / RSA Encryption}

We have:

\begin{align*}
    M: &\; \text{Message} \\
    C: &\; \text{Encrypted Message} \\
    e: &\; \text{Encryption Key} \\
    d: &\; \text{Decryption Key} \\
    p: &\; \text{Prime mod 1} \\
    q: &\; \text{Prime mod 2} \\
    N: &\; \text{Product of } p \cdot q \\
\end{align*}

To \textbf{encrypt} a message $M$ to $C$

$$C \equiv M^e \pmod{N}$$

To \textbf{decrypt} $C$ back into $M$, we must first find $d$

$$ed \equiv 1 \pmod{(p - 1)(q - 1)}$$

Just find $e^{-1}$ using Euclidean and plug it into

$$M \equiv C^d \pmod{N}$$

\subsection{Proof By Induction}

$$\sum_{k = 1}^{n}k^3 = \left[ \frac{n (n + 1)}{2}\right]^2$$

\begin{poof}
    BASE CASE:

    \begin{align*}
        \sum_{k = 1}^{1}k^3 &= 1^3 = 1 \\
        &= \left[ \frac{n (n + 1)}{2}\right]^2 = \left[ \frac{1(1 + 1)}{2}\right]^2 = 1 \\
    \end{align*}

    So the case $n = 1$ is true

    INDUCITVE CASE:

    Assume that $\sum_{k = 1}^{n}k^3 = \left[ \frac{n (n + 1)}{2}\right]^2$

    We then prove $\sum_{k = 1}^{n + 1}k^3 = \left[ \frac{n + 1 (n + 2)}{2}\right]^2$

    Let's do this by making the more complicated part (the proving part) equal to the simpler one

    \begin{align*}
        \sum_{k = 1}^{n + 1}k^3 = \sum_{k = 1}^{n}k^3 + (n + 1)^3 &= 
        \sum_{k = 1}^{n + 1}k^3 = \left[ \frac{n(n + 2)}{2}\right]^2 + (n + 1)^3 \\
        &= \frac{n^2(n + 1)^2}{4} + (n + 1)^3 \\
        &= (n + 1)^2 \cdot \left[ \frac{n^2}{4} + (n + 1)\right] \\
        &= (n + 1)^2 \cdot \left[ \frac{n^2 + 4n + 4}{4} \right] \\
        &= (n + 1)^2 \cdot \frac{(n + 2)^2}{4} \\
        &= \frac{(n + 1)^2(n + 2)^2}{4} \\
        \Aboxed{&= \left[\frac{(n + 1)(n + 2)}{2}\right]^2}
    \end{align*}
\end{poof}

\subsection{Multiples of n up to k}

\begin{theorem}[Number of Multiples of n up to k]
    To find the number of positive integers divisible by $n$ up to $k$ is given by

    $$\left\lfloor \frac{k}{n} \right\rfloor$$

    This DOES NOT include 0, so for all whole numbers it would be

    $$\left\lfloor \frac{k}{n} \right\rfloor + 1$$
\end{theorem}

\subsection{Pigeonhole Principle}

\begin{definition}[Pigeonhole Principle]
    If there are $n$ boxes and more than $n$ objects places in these boxes, at least one box contains more than one object.

    At least one box will contain 

    $$\lceil \frac{k}{n} \rceil$$

    objects
\end{definition}

\subsection{Permutaions and Combinations}

\subsubsection{Permutations}

\begin{definition}[Permutation]
    A \yellowbox{Permutation} is one specific way a group of items can be arranged where order matters. 
    In general, if we have $n$ objects and we want to arrange $k$ of them, then the number of ways is given by

    $$\prescript{}{n}{P}_k = \frac{n!}{(n - k)!}$$
\end{definition}

\subsubsection{Combinations}

\begin{definition}[Combination]
    A \yellowbox{Combinaton} is just an arrangement where order does not matter.

    The number of combinations of $k$ items taken from $n$ items is

    $${n \choose k} = \prescript{}{n}{C}_k = \frac{n!}{k!(n - k)!}$$

    You would say "n choose k"

    Think of this as "how many ways can you arrange $k$ things in $n$ slots"
\end{definition}

\subsection{Binomial Theorem}

\begin{theorem}[Binomial Theorem]
    
    $$\text{For } n = 0, 1, 2, \ldots, (a + b)^n = \sum_{k = 0}^{n} {n \choose k} a^{n - k}b^k$$

\end{theorem}

\subsection{Sum of Combinations}

\begin{theorem}[Sum of Combinations]
    $$\text{For } n = 0, 1, 2, \ldots, \sum_{k = 0}^{n}{n \choose k} = 2^n$$
\end{theorem}

\subsection{Arrangements of Letters}

For a word with $k$ letters, where each letter has $n_i$ appearances, the numbers of arrangements is given by

$$ \frac{k!}{n_1! \cdot n_2! \cdot \ldots \cdot n_k!}$$

You basically just do

$$ {k \choose n_1} \cdot {k - n_1 \choose n_2} \cdot {k - n_1 - n_2 \choose n_3} \cdot \ldots$$

For a word with all unique letters, it's just

$$k!$$

\subsection{Stars and Bars}

When you have $n$ slots and $k$ options for each slot, the number of possible combinations is 

$$\frac{(n + k - 1)!}{n! \cdot (k - 1)!} = {n + k - 1 \choose k - 1}$$

Another way of thinking of this is stars and bars

Let

\begin{align*}
    s: &\; \text{Stars, num slots} \\
    b: &\; \text{Bars, options - 1} \\
\end{align*}

Then the combinations is

$$\frac{(s + b)!}{s! \cdot b!} = {s + b \choose b}$$

\subsection{Binomial Probability}

Given a limited event (ie 2 outcomes) with probability $p$,

$$p(k \text{ successes in } n \text{ trials}) = {n \choose k} \cdot p^k \cdot (1 - p)^{n - k}$$

\subsection{Expected Value (EV)}

$$EV = \sum (\text{probability}) (\text{payoff})$$

% ========== End the document ==========
\end{document}